\section{the direct strategy closed, and the fence posts}\label{sec:fenceposts}

\begin{Lem}{13}[The spectral-moment method is the wrong invariant]
Studying the multiplication operator $L_u\colon H\mapsto B(u,H)$ gives $\tr(L_W)=0$, $\tr(L_W^2)=\tfrac{(n-1)(2n-1)}4\|W\|^2$, and a Peirce cap at $\lambda/2$ on $u^\perp$ attained on a subspace of dimension exactly $\dim(O(n)\cdot u)$. Together these yield $\lambda^2(1+\tfrac{k_0}4)\le\tfrac{(n-1)(2n-1)}4$ from criticality alone. Since $k_0\le\binom n2$, the method's ceiling is $\lambda\le2-o(1)$ against a true $\sqrt2$, an excess of $58.9\times$ at $n=12$ and $3.09\times$ asymptotically, with the fourth moment reaching only $14\times$ even granting an oracle.
\end{Lem}

\begin{proof}
Let $\lambda=P(u)$ at a critical point $u$ and let $k_0=\dim\big(O(n)\cdot u\big)$. The three ingredients are: $u$ is itself an eigenvector of $L_u$ with eigenvalue $\lambda$; the Peirce cap gives $\langle L_uH,H\rangle\le\tfrac\lambda2\|H\|^2$ for every $H\perp u$ at a local maximum, and it is \emph{attained} on a subspace of dimension exactly $k_0$; and $\tr(L_u^2)$ is the fixed constant above. Summing squares of eigenvalues over the $\lambda$-eigenvector and the $k_0$-dimensional cap eigenspace,
\[
\lambda^2+\tfrac{\lambda^2}4k_0\;\le\;\tr(L_u^2)=\tfrac{(n-1)(2n-1)}4 ,
\]
which is the stated inequality. Since $k_0\le\binom n2$, the best conclusion the method can ever reach is $\lambda\le2-o(1)$, against the true value $\sqrt2$.
\end{proof}

The diagnosis is structural rather than a matter of tuning: $\tr(L_W^2)$ is \emph{constant} in $W$ on the unit sphere, so it carries no information beyond the norm, and no refinement of a method built on it can distinguish one unit Weyl operator from another. The fourth moment reaches only $14\times$ even granting an oracle for the unknown quantities.

\paragraph{The third moment closes the route completely, and the closure is arithmetic.} $\tr(L_x^3)$ is determined by the two cubic invariants, so the moment data available is exhausted at order three: a fourth universal identity does not exist. The constraint the third moment places on $b=\langle W^{\#},W\rangle$ is $b\le 2n\theta/(n-1)+4\theta^3/(n^2-1)$, and it is \emph{dominated} by one line: $\sum\lambda^2=\|W\|^2=1$ forces $|a|\le1$, hence $b=\theta-a\le\theta+1$, which is smaller at every $n\ge6$ ($2.24$ against $2.75$ at $n=12$). \textbf{So the moment hierarchy adds nothing after the second moment, and its remaining orders are unavailable rather than merely weak.}

\paragraph{The extremal-idempotent calculus does not extend here either.} Euler--Jacobi requires finitely many simple roots and no nonzero $2$-nilpotents, and \emph{both} hypotheses fail on $\Weyl_n$: the $O(n)$-orbit of an idempotent is a positive-dimensional family of idempotents, so the root set is infinite; and $\Weyl_n^{\mathbb C}$ carries nonzero $2$-nilpotents, since complex orbits of idempotents are positive-dimensional hence unbounded. \emph{No equivariant form of the identity exists in the literature: global-residue extensions handle isolated degenerate zeros and never positive-dimensional ones.} Lemma~24 continues to arrive from \cite{kt2018} as a pointwise eigenvalue bound, which is the part of that calculus this problem can use.

\begin{remarkstar}
One genuine by-product survives the closure of this route, and it is a theorem about the problem rather than about the method. The cubic $O(n)$-invariants on $\Weyl_n$ are spanned by $\tr(W^3)$ and $\langle W^{\#},W\rangle$ in \emph{every} dimension. Only their rank changes, being $1$ at $n=4,5$ (where $\langle W^{\#},W\rangle=2\tr(W^3)$ identically) and $2$ for $n\ge6$. The proof combines Weyl's first fundamental theorem (which reduces every cubic invariant to one of $1728$ complete contractions of $W^{\otimes3}$, in $5$ orbits under $S_3\times D_4^3$), monotonicity of the invariant space under zero-extension, the Newell--Littlewood count $N_{(2,2),(2,2),(2,2)}=2$ in the stable range, and exact integer arithmetic on the full symmetric polarization at $n=4,5,6$.
\end{remarkstar}

\subsection{Why the one sharp bound in the literature stops at \texorpdfstring{$n=5$}{n=5}}\label{sec:collapse}

Exactly one sharp constant $\thstar$ exists in the consulted literature: Huisken's Lemma 3.5 \cite{huisken1985}, proved on $\Weyl_4$. It is worth stating why it does not travel, and how far it does. \textbf{The ceiling is $n=5$, not $n=4$}, and the mechanism is a dimension count on an antisymmetriser (Lemma~51) rather than the chirality splitting of $\so(4)$, which explains $n=4$ and cannot reach $n=5$.

\paragraph{At $n=4$ the cubic is spectral.} $\so(4)=\so(3)\oplus\so(3)$ splits $W$ into chiralities $W^{\pm}$, each a symmetric \emph{trace-free} operator on $\so(3)$, and on a $3$-dimensional ideal the structure constants are a multiple of the volume form, so $\#$ restricts to a multiple of the adjugate, concretely $c^{\pm}=\mp\sqrt2\,\varepsilon$ and $\#=2\operatorname{adj}$ on each chirality. Since $\langle\operatorname{adj}A,A\rangle=3\det A$ and $\tr(A^3)=3\det A$ for trace-free $3\times3$ $A$,
\[
\langle W^{\#},W\rangle=6\big(\det W^{+}+\det W^{-}\big)=2\tr(W^3),\qquad\text{so}\qquad \langle Q(W),W\rangle=3\tr(W^3),
\]
a purely spectral quantity. Lemma 3.5 then has a two-line proof (Okumura's sharp bound at $m=3$ on each chirality, then $a^{3/2}+b^{3/2}\le(a+b)^{3/2}$), giving $\thstar$ with the $\CP^2$ equality case for free, and Huisken's own enumeration of stationary points by Lagrange multipliers is not load-bearing. \emph{A norm caution}: his $|W|$ is the full tensor norm and this report's $\|W\|$ the operator norm on $\Lambda^2$, with $|W|=2\|W\|$, so his $\sqrt6/4\,|W|^3$ is $\thstar\|W\|^3$ here.

\begin{Lem}{49}[The collapse space, and its dimension]
Write $\theta=\sum_\alpha\lambda_\alpha^3+3\sum_{\alpha<\beta<\gamma}T_{\alpha\beta\gamma}\lambda_\alpha\lambda_\beta\lambda_\gamma$ in an eigenframe, with $T_{\alpha\beta\gamma}=\langle[v_\alpha,v_\beta],v_\gamma\rangle^2$. The condition that $\theta$ reduce to a multiple of $\sum_\alpha\lambda_\alpha^3$ is \emph{linear in $T$}, so the collapsing configurations form a subspace, and its dimension is governed by the number of independent trace conditions the spectrum obeys, which is the number of irreducible summands of $\Weyl_n$. Under a single such condition the collapse space is exactly one-dimensional, $T$ constant, at every $N\ge4$, and constant $T$ is not realised by any $\so(n)$.
\end{Lem}

\begin{proof}
A Weyl tensor is trace-free, so $p_3=3e_3$ modulo $(e_1)$ with both sides squarefree, and the question is for which $T$ the squarefree cubic $\sum_{\alpha<\beta<\gamma}T_{\alpha\beta\gamma}\lambda_\alpha\lambda_\beta\lambda_\gamma$ equals $\kappa\,e_3$ on the relevant $\lambda$-space. That is a linear system in the $\binom N3$ unknowns $T_{\alpha\beta\gamma}$, solved exactly below.

Constant $T$ is unrealisable. The row marginal $\sum_{\beta\gamma}T_{\alpha\beta\gamma}=2(n-2)$ would force $t=\tfrac{16}{(n+1)(n^2-n-4)}$, but $T_{\alpha\beta\gamma}$ vanishes on every triple whose bracket is orthogonal to $v_\gamma$, and in $\so(n)$ most triples are of that kind: the standard basis gives $T\in\{0,1\}$, and $\su(3)$ in the Gell-Mann basis gives $T\in\{0,\tfrac14,\tfrac34,1\}$ with $47$ of its $56$ triples at $0$.
\end{proof}

\begin{certificate}
Exact linear algebra over $\mathbb Q$ on the $\binom N3$-dimensional $T$-space.
\begin{center}
\begin{tabular}{@{}ll@{}}
\toprule
$\lambda$-space & collapse dimension of $\binom N3$ \\
\midrule
$N=6$, one global trace condition & $\mathbf1$ of $20$ \\
$N=6$, two conditions $(3+3)$ (the $\so(4)$ chirality case) & $7$ of $20$ \\
$N=6$, three conditions $(2+2+2)$ & $14$ of $20$ \\
$N=7$, one global & $\mathbf1$ of $35$ \\
$N=7$, two $(3+4)$ & $10$ of $35$ \\
$N=9$, three $(3+3+3)$ & $38$ of $84$ \\
\bottomrule
\end{tabular}
\end{center}
\end{certificate}

\paragraph{What Lemma~49 settles, and what it does not.} $\Weyl_4=\mathcal W^{+}\oplus\mathcal W^{-}$ carries \emph{two} independent trace conditions, giving the $7$-dimensional collapse space that the chirality splitting lives in, whereas $\Weyl_n$ is $O(n)$-irreducible for every $n\ge5$, so there is exactly one condition and the collapse space is the single line $T\equiv\text{const}$, which $\so(n)$ does not realise. \textbf{That is a statement about $\lambda$-uniform mechanisms and it decides no dimension.} Collapse needs $\sum T\lambda\lambda\lambda=\kappa e_3$ only at the single $\lambda(W)$ that $T(W)$ comes with, and Lemma~51 shows the cubic does collapse at $n=5$, where $\Weyl_5$ is $O(5)$-irreducible and the realised $T$ is manifestly non-constant. So irreducibility is not the separating fact, but the rank of the six-index Kronecker delta is (\S\ref{sec:lovelock}). \textbf{What is true in both readings is that the mechanism dies at $n=6$}, and that an extension past it would have to be an \emph{approximate} collapse ($T$ near constant, giving a bound rather than an identity). The $n=6$ isospectral family of post~1 measures the distance to be covered: one spectrum, and $P$ ranging over $\{-24,0,+24\}$. \emph{The second spectrum is what covers it}, and \S\ref{sec:bispectralid} is that statement.

\medskip
The following proven or machine-verified negatives delimit what any new attempt must clear. There are thirty-one, and a route that one of them already kills is the most expensive kind of proposal, because it looks new.

Each post states the negative that holds, and where a natural stronger reading of it is false, the post says so in place: posts 17, 18 and 19 hold with reasoning or hypotheses different from the ones they invite, and post F20 holds only in a narrower form than its name suggests. \textbf{A post is kept at its number even when it is the narrow form that survives}, because a negative that disappears from the list is paid for a second time.

\begin{enumerate}[leftmargin=1.8em]
\item \textbf{A bound reaching $P$ only through $\tr(W^3)$ cannot close the gap.} At the conjectured maximizer essentially all of $P$'s value sits in $\langle W^{\#},W\rangle$, the structure-constant term, which depends on eigenspace \emph{alignment} rather than on the spectrum, while the Tachibana/H\"older estimates bound $\tr(W^3)$ alone, and there $\tr(W^3)/P=\tfrac2{(n-1)(n-2)}$, an identity rather than a scan. \emph{The witness is exhibited rather than argued}: at $n=6$ there is a family of $300$ diagonal Weyl tensors sharing the $\Lambda^2$-spectrum $(1^6,0^3,(-1)^6)$, all with $\|W\|^2=12$, $\tr(W^3)=0$ and trivial stabiliser, on which $P$ takes the values $-24$, $0$, $+24$ with counts $60/180/60$. So \textbf{$P$ is provably not a function of the spectrum}. \emph{Scope, and it is narrow}: that says no spectral \emph{function} determines $P$, but a spectral \emph{bound} can still be sharp, and at $n=4$ one is (Lemma~49). The post says nothing about bounds on $\langle Q(W),W\rangle$ taken whole. Post F23 is that route.
\item \textbf{The two separately-sharp cubic bounds overshoot by $1.7$--$1.8\times$} at all $95$ integers $n\in[6,100]$, \emph{and all of the slack sits in the $\tr(W^3)$ bound}. At the balanced sphere product the vertex bound is exactly saturated:
\[
\frac{B_{\#}}{\theta_0}=\frac{n(n-3)}{(n-1)(n-2)}=1-\frac2{(n-1)(n-2)},\qquad \frac{\tr(W^3)}P=\frac2{(n-1)(n-2)},
\]
the two summing to $1$ identically. The $\tr(W^3)$ bound is sharp for arbitrary trace-free operators and loose on Weyl tensors by at most $2.32\times$ at $n=8$ and $2.39\times$ at $n=12$. Using only \emph{attained} values, any term-by-term route overshoots by $f(n)=\tfrac{n(n-3)}{(n-1)(n-2)}+\tfrac{n\sqrt3}{6\sqrt{(n-1)(n-2)}}$, which is unimodal, peaks at $n=11$, and decreases to the \emph{limit} $1+\tfrac{\sqrt3}6=1.2886\ldots$ from above (exactly, $f(323)>\tfrac{129}{100}>f(324)$). It is a limit and not a floor: $f(4)$, $f(5)$ and $f(6)=\tfrac{9+\sqrt{15}}{10}$ lie strictly below it, and $n=6$ is inside the live range, so the figure applies from $n\ge7$. \textbf{A proof needs a joint sharp bound on $\langle Q(W),W\rangle$, not a term-by-term sum}, and hunting compatible extremal profiles is the wrong repair, the profiles already coinciding.
\item \textbf{Universal lifting is false}, proven three independent ways, and global polarity fails by dimension count already at $n=4$.
\item \textbf{{\L}ojasiewicz constants are astronomically large}, so that route cannot produce a usable separation.
\item \textbf{Chain-tightening caps at $\delta^2\lesssim10^{-3}$}, short of what the defect-penalty route required.
\item \textbf{Spectral moments of $L_W$ are the wrong invariants} (Lemma~13), $58.9\times$ short at $n=12$.
\item \textbf{Every candidate inequality must see the crossover $n=6+2\sqrt7$.} An inequality that does not change character there is bounding the wrong thing, since below the crossover the rival genuinely wins.
\item \textbf{The sharp (II)-side bound is Fox-sharp but does not separate}: $(S^m)^3$ reaches $1.234230>\thstar$ at $n=24$. Any repair must be $n$- or value-aware.
\item \textbf{Marginal relaxations of the bracket tensor cannot bound $P$.} Writing $P$ in the eigenframe through $T_{abc}=c_{abc}^2$ and imposing every constraint the true tensor satisfies (nonnegativity, total symmetry, row marginals $2(n-2)$, pair caps $2$, and $\tr(W^{\#})=-1$) still leaves a linear-programming optimum above $\theta_0$ by a factor near $0.80\sqrt{n-2}$. What the constraints omit is frame-realizability, which is the original problem. A route that relaxes alignment into marginals is already dead.
\item \textbf{Enumeration cannot close the problem: a bound is needed, not a census.} The critical set is exponentially large and mostly symmetry-free: $20$ unpenalized restarts at $n=6$ gave $17$ distinct critical values, every one with trivial stabilizer. And spectral structure is not a reliable signal of group-theoretic origin: a zero-defect critical point was found with three distinct eigenvalues of multiplicities $5,6,4$ and yet a trivial stabilizer.

\item \textbf{Splitting $P$ into its two terms cannot bound $M(n)$ below $1+\sqrt2$}, where $M(n)=\max\{P(A):A\in\mathcal W_n,\ \lVert A\rVert=1\}$ is the global maximum, positive at every $n\ge4$ and distinct from the diagonal maximum $\Theta_n$. At unit norm $\tr W^3\le1$, but the $\#$ term carries almost all of $P$ where $P$ is large ($\langle W,W^{\#}\rangle=1.2136$ against $\tr W^3=0.0225$ at $S^6\times S^6$), and it tends to $\sqrt2$ as $n\to\infty$. Any route to $\sqrt3$ must treat the two terms together.

\item \textbf{A single flatness inequality cannot cap $M(n)$.} The threshold $y_*(n)$ solving $\sqrt y+(n-2)y=\sqrt3$ is $0.1363$ at $n=12$ against a floor of $0.1231$, and \emph{six of the eight} catalogued critical points exceed it. The flat and low-rank regimes need separate arguments.

\item \textbf{Multistart root-finding cannot census the diagonal $S=0$ set, and the reason is structural.} Degenerate critical points have singular Jacobians, so Newton never lands on them: the $C_5$ blow-up is diagonal, exactly critical, triangle-free and block-constant, and an honest census at $n=10$ must return it. It never does. Likewise $1200$ starts on $\Weyl_6$ reached a maximum of $0.408$ with $\CP^3$ (an $n=6$ critical point at $\theta^2=10/9$ with $S=2\times10^{-30}$) never appearing. \textbf{Random-start Newton does not reach the high-$\theta$ basin, so such an enumeration is evidence of completeness in neither direction.}

\item \textbf{Dropping Ricci-flatness is safe for the bound and fatal for $\theta$.} The same maximisation over the curvature-operator subspace instead of $\Weyl_n$ returns $\sqrt{2(n-1)/n}$ with $S=0$ exactly: the \emph{round sphere}, $W=I/\sqrt N$, which has no negative eigenvalues and is not Ricci-flat. Since $\sqrt{2(n-1)/n}>\theta_0(n)$, relaxing to Bianchi alone would manufacture a refutation of the conjecture out of a tensor that is not a Weyl tensor.

\item \textbf{$\lambda_{\max}(M)$ cannot be certified from $T$'s marginals.} Adding the full vector identity $R\lambda=-2\lambda$ to the transportation relaxation ($N$ linear equations, of which $\tr(W^{\#})=-1$ is the single $\lambda$-contraction) moves the $n=6$ optimum from $1.722739$ to $1.413614$, an $18\%$ cut, and \textbf{the tightening does not scale}: at $n=8$ it is $1.4\%$ while the overshoot grows from $1.63\,\theta_0$ to $2.00\,\theta_0$. At the padded $W_{\CP^2}$ spectrum the relaxation returns exactly $\thstar$ under both constraint sets, so the whole overshoot sits at spectra that are not Weyl spectra. Frame-realizability has to enter.

\item \textbf{The two-negative full-support stratum is measure zero, so no sampling method can census it.} With the two negative eigenvectors fixed, existence is an exact convex feasibility problem, and random draws over every rank profile at $m=6$ return $0/36$ feasible, \emph{including the profile where a witness demonstrably exists}. Objects in this stratum must be constructed, and any negative result from a search over it is a statement about the search.

\item \textbf{The three per-index relations on the bracket tensor are complete, and no fourth exists.} Writing $P$ in the eigenframe through $T_{abc}=c_{abc}^2$, the relations are Killing row mass $\sum_{bc}T_{abc}=2(n-2)$, the per-row identity $\sum_{bc}\lambda_bT_{abc}=-2\lambda_a$, and criticality $\sum_{bc}\lambda_b\lambda_cT_{abc}=2\lambda_a(\theta-\lambda_a)$. Together they pin every contraction $\sum\lambda_a^p\lambda_b^q\lambda_c^r T_{abc}$ with two exponents at most $1$, and the four hunting grounds for a fourth (Bianchi, Jacobi's square-producing contractions, the off-diagonal criticality block, and the quartic Casimir) are each closed, the last three by termwise vanishing, by averaging to criticality, and by circularity through $Q$ preserving $\Weyl_n$. \emph{The consequence for the route is a ceiling rather than an inertness}: the full pinned relaxation reaches $1.5$--$3.7\times\thstar$, so \textbf{a successor must use a $\{2,2,0\}$ moment or higher}, and naming its lever's weight pattern decides the matter in advance.

\item \textbf{Only the \emph{unweighted} generator sum is inert, and it needs $n\ge5$.} Summing Petersen--Wink's $\lVert[\mathrm{ad}(v_\gamma),W]\rVert^2$ over an orthonormal basis of $\so(n)$ returns the Casimir constant $4(n-1)$, which is not new information: it follows from the Killing-form row marginal $2(n-2)$ and $A(W)=2W$. The scoping is load-bearing at both ends. \emph{The hypothesis}: at $n=4$ the adjoint of $\so(n)$ is reducible and tracing separately over $\Lambda^2_\pm$ gives different answers while the total stays pinned. \emph{The scope}: the very next weight is not inert at all, since $\sum_\gamma\lambda_\gamma g_\gamma=-4\theta$ returns $\theta$ itself, which is Lemma~41.

\item \textbf{Integrating $\theta$ over a manifold gives a lower bound, and the conjecture is pointwise-algebraic}, so the Weitzenb\"ock route is closed. \emph{The sign recorded here was inverted and the hypothesis was too weak.} With the corrected constant $\Theta_3=-2\theta$, Weingart's Lemma 6.3 gives on a compact \emph{Ricci-flat} (not merely parallel-Ricci) manifold $\int\theta=\tfrac12\lVert\nabla R\rVert^2\ge0$. The route stays closed for the opposite reason to the one recorded.

\item \textbf{Only the \emph{uniform single-plane} third moment fails to determine $\theta$.} The wider reading, that $\theta$ is not a closed form in sectional-curvature moments at all, is false, because a second independent point of the cubic $2$-plane exists: the three-plane correlation
\[
\Phi=\mathbb E_{x,y,z}\big[\sec(x,y)\,\sec(y,z)\,\sec(z,x)\big],
\]
which gives $\theta=c_n\Psi_3+\tfrac5{72}\Phi$ exactly. What the post does say is sharp: $\Psi_3$ alone does not determine $\theta$, since $S^3\times S^3$ and $\CP^3$ share $\theta^2=\tfrac{10}9$ \emph{exactly} while their $\Psi_3$ differ by a factor of \emph{exactly} $13.5$, computed by Wick sums with no sampling. \textbf{The wide form is stated here because it is what an attempt on this route reaches for first.}

\item \textbf{Sectional-curvature machinery cannot deliver $|\lambda_{\min}(W)|$.} The sharp Weitzenb\"ock/PIC2 literature is the natural source, and the identification is real: $q_2(W)=-2W$ on Weyl tensors, verified at $10^{-16}$, so their operator is ours. But every estimate in that family is phrased in \emph{sectional} curvature, and at the K\"ahler points $\lambda_{\min}$ is attained on the \textbf{non-decomposable} K\"ahler-form bivector that no $2$-plane witnesses: $K_{\min}=-0.079$ against $\lambda_{\min}=-0.211$ at $\CP^3$, and $-0.042$ against $-0.141$ at $\CP^4$. \textbf{So the machinery is blind at exactly the points that set the $\thstar$ threshold}, and it is quantitatively $2.5$--$7.8\times$ short besides.

\item \textbf{$\ker\Phi_W$ decides diagonality and carries no information about $\theta$.} Lemma~37's operator is the first frame-independent grip the programme has had on diagonality, and it was named as the handle for \emph{characterising} the non-diagonal $S=0$ set. Measured, it does not track $\theta$: $\dim\ker\Phi_W=1$ is the generic value, holding at every $\CP^m$ \emph{and} at random Weyl tensors whose $\theta$ differ by an order of magnitude, and since $\Id\in\ker\Phi_W$ always, that value forces $\ker\Phi_W=\R\cdot\Id$ exactly, so no refinement of the kernel can separate them. \textbf{Use it to define the target set, not to bound the functional on it.}

\item \textbf{The published bound on $\langle Q(W),W\rangle$ taken whole owes a factor growing like $n/\sqrt2$, not a constant.} Xu's Lemma 2.14 never splits the cubic and gives $|\langle Q(W),W\rangle|\le\sqrt{(n^2-1)(n-2)/n}\,\|W\|^3$, overshooting $\theta_0$ by $\sqrt{n(n+1)/2}$ ($\sqrt{10},\sqrt{21},\sqrt{78},\sqrt{210}$ at $n=4,6,12,20$). Posts 1 and 2 do not cover it, since they concern the split. The whole loss is localised: his (2.53) and (2.54) are identities, and (2.55) applies $\sqrt{(N-1)/N}\,\|W\|$ to every $\xi^{pqrs}$ independently, as if each could simultaneously be $W$'s top eigenvector. Writing $M=\sum\xi\otimes\xi$, a positive semidefinite operator with $\tr M=16(n-1)\|W\|^2$ and $\tr(WM)=-16\langle Q(W),W\rangle$, the conjecture reads exactly $|\tr(W\widehat M)|\le\theta_0(n)\|W\|/(n-1)\sim\sqrt2\|W\|/n$ with $\widehat M=M/\tr M$, against the $\approx\|W\|$ that (2.55) supplies. Needed over used is $0.2108$, $0.1124$, $0.0688$, $0.0280$ at $n=6,12,20,50$. \textbf{Closing it means using that $M$ is quadratic in $W$}, which is the alignment information (2.55) discards. As a statement about $\theta$ the $M$-form is Conjecture~A rewritten, neither easier nor harder. What it supplies is the measurement.

\item \textbf{The collapse of $P$ to a spectral quantity is confined to $n\le5$, and no classification of the structure constants changes that.} \emph{The dimension was corrected from $n=4$ to $n\le5$ at Gate~2, and the driver with it.} Lemma~51 gives the true mechanism: $\beta-\tfrac\alpha2$ is a Lovelock dimensionally-dependent identity whose six-index Kronecker delta vanishes identically for $n\le5$, so the cubic is a function of the $\Lambda^2$-spectrum alone there and at no larger $n$. Lemma~49's collapse-space computation is right (a subspace of $T$-space whose dimension is the number of irreducible summands of $\Weyl_n$, two at $n=4$ and one at every $n\ge5$), but it counts $\lambda$-\emph{uniform} mechanisms and therefore decides no dimension: $\Weyl_5$ is $O(5)$-irreducible, its realised $T$ is non-constant, and it collapses regardless. \textbf{The post is kept because the negative it names holds}, one dimension further out than it said. \emph{And the missing ingredient is not an invariant theory of the bracket $3$-form}: there is exactly one such form per $n$, known in closed form, so nothing is being classified; the forms obey the Jacobi identity and so sit inside the Cartan--Killing classification, finite in every dimension; the relevant group is $\mathrm{Aut}(\so(n))\subset O(N)$, an eigenframe change, not $GL_N$; and \textbf{moduli do not preclude bounds}, since $n=4$ carries $\tfrac{n^2(n^2-13)}{12}=4$ moduli \emph{and} a proved sharp bound. The moduli count $\tfrac{n^2(n^2-13)}{12}$ is correct arithmetic ($4,25,69,272,1572$ at $n=4,5,6,8,12$), and orbits at the maximizers are \emph{smaller} than generic ($n^2/4$ and $4n-14$ against $n(n-1)/2$), so their slices are larger, $1602$ and $1604$ against $1572$ at $n=12$.

\item \textbf{A sample that cannot contain the target is not evidence about the target.} Four long runs at $n=6,8,10,12$ screened $2{,}060{,}095$ Weyl tensors for a falsifier of $\sup\{\theta:Q(W)=\theta W,\ S=0,\ W\text{ non-diagonal}\}\le\thstar$, and \emph{zero} lay in the target set: none was critical, the sampler taking $60$ fixed-step gradient steps and finishing at residual $\approx4\times10^{-2}$ against a $10^{-8}$ tolerance, and none had $S=0$, which is measure zero for random draws in any case. So the runs could not have produced a falsifier however long they ran, and the recorded maximum pressing to $10^{-5}$ below $\thstar$ at $n=6$ is the global supremum over $\Weyl_6$, already known, and says nothing about $\{S=0\}$. \textbf{Before a sample is quoted as evidence, check that the sample can contain the thing being looked for}. A screen applied afterwards does not rescue a population that never met the hypothesis. Here the screen was present and refused every point. What was missing was noticing that it refused \emph{all} of them.

\item \textbf{Improving the second-kind side alone cannot close the conjecture, and a bound exact at the coordinate-diagonal end cannot either.} Both negatives come from the bi-spectral form $\theta=\tfrac32p_3(\mu)-\tfrac12p_3(\omega)$ (Lemma~50) and both are provable rather than sampled. \emph{One-sided}: $\tfrac32\max p_3(\mu)=1.4657$ already exceeds $\theta_0(12)=1.236033$, so however sharp an $\omega$-side estimate becomes, the $\mu$ term alone overshoots. \emph{Exact-at-the-end}: writing $t=\|K\|_F^2\le2$ for the vertex matrix, $t=2$ holds exactly when $W$ is coordinate-diagonal, the Schur--Horn relaxation is forced to $\theta$ there by equal sums and equal norms (Lemma~54), and it is non-increasing in $t$ because the vertex class is a linear space. \textbf{So any bound that is exact at $t=2$ and monotone in diagonal information is $\ge\theta_0(n)$ at every $t<2$} (three lines, and nothing in them is specific to Schur--Horn or to Okumura). \emph{Two corollaries worth naming separately}: a \emph{lower} bound on the frame-optimized diagonal spread is worth exactly nothing, the relaxation staying at $4.0129$ for every constraint level up to the hard ceiling; and monotonicity in $t$ is a theorem, so passing it is not a validity gate, because it tolerates under-exploration exactly, which is the failure mode it met.

\item \textbf{No bound of the form $\theta\le\Pi/4$ closes the target}, with $\Pi=|\lambda_{\min}|(4(n-1)-\Delta-G_0)$ built from the two contractions $\sum g=4(n-1)$ and $\sum\lambda g=-4\theta$. The obstruction is the reverse of looseness: the slack is exactly
\[
\Pi-4\theta=\sum_{\mm}\big(|\lambda_{\min}|-|\lambda_\gamma|\big)g_\gamma+\sum_{\mathfrak p}\lambda_\gamma g_\gamma,
\]
both terms non-negative, vanishing exactly when $\Delta=0$ and $g$ is carried by the $\lambda_{\min}$ eigenspace. \textbf{Those are the graded points}: every eigenvector outside that eigenspace is then an infinitesimal symmetry. At $\CP^m$ the symmetries number $m^2$ against multiplicity $m(m-1)$, at $S^k\times S^k$ they number $\dim(\so(k)\oplus\so(k))$ against $k^2$, in both cases the isotropy and the tangent space of the symmetric pair. So $\Pi/4=\theta$ identically on the graded half, measured slack $\le1.7\times10^{-15}$ across the catalogue, and there the bound is the target \emph{restated} rather than relaxed. \textbf{It is strictly slack only on the ungraded half}, which is (G). \emph{A separate, weaker negative}: $\Pi$ is outright unbounded on $\{S=0$, non-diagonal, critical$\}$ via $\perp$-sums, criticality forcing $1/\theta^2=\sum1/\theta_i^2$, but those points cannot be maximizers, so that is not what closes the route.

\item \textbf{A uniform-in-$n$ spectral gap on the graded critical set does not exist, and the conjectured maximizer is the witness.} Closedness of the graded set reduces to a lower bound on the smallest \emph{nonzero} eigenvalue modulus over that set (nonzero because padding a graded critical point leaves it graded and contributes exact zeros, so the raw infimum is $0$ and the uniform statement is vacuous). On the nonzero spectrum the reduction survives padding, and there the balanced sphere product supplies the shrinking sequence: its smallest nonzero modulus is $\sqrt2(n-2)/\big(n\sqrt{(n-1)(n-2)}\big)\sim\sqrt2/n\to0$, exactly at $n=6,8,10,12,16$. \textbf{Any usable form of the statement is per~$n$.} \emph{Whether a per-$n$ floor read off the named families describes the whole graded set is open, and is blocked on method rather than evidence}: every graded point at $n=6$ is a \emph{degenerate} zero of the criticality equation (excess $9$, against $0$ at every ungraded point found), so a residual-cutoff root-finder cannot reach one at any compute, and gradient ascent cannot either, $S^3\times S^3$ being a degenerate saddle (\S\ref{sec:closedness}, and post F13 rescoped).

\item \textbf{No mass on $\pp$ separates the pair-graded critical points from the merely sign-graded ones.} $\|W|_{\pp}\|^2$ is not a function of the point: the $n=8$ witness admits at least eight inequivalent $\so(3)\oplus\so(5)$ splittings, and the value spans $[0.319349,\,0.808897]$ across them, straddling the one the equivariant point carries. At the equivariant point itself the value is $(p^2+q^2-n)/((n-1)(n-2))$ exactly (confirmed at eleven pairs from $n=6$ to $14$), and that is a Cauchy--Schwarz floor, saturated because Schur makes $W$ scalar on each simple factor. The floor carries no information off that locus, the block traces not being invariants of the pair. \emph{A successor separates the two by a quantity that is not a mass, and names its splitting.}

\item \textbf{The $\pp$-internal cross-block is not a usable test direction.} It vanishes by Schur exactly at equivariance, which makes it look like the structural quantity the previous post asks for. At the $n=8$ witness $57.7\%$ of its mass lies outside $\Weyl_n$, and after projecting back and removing the $W$-component its Rayleigh quotient is $-0.290085$, against the $\tfrac12$ a certificate needs and the $0.608885$ attained there. \emph{On this problem $\Delta>0$ enters as a hypothesis about the isotropy and never as a supply of test directions (the sixth route to close for that reason).}

\item \textbf{Lemma~23's bound does not extend off the two-eigenvalue family in the form $\theta\le\frac2n\sqrt{rN/(N-r)}$ with $r=\dim E_+$.} The obstruction is padding: $\theta$ and $\dim E_+$ are padding-invariant while the bound behaves like $\frac2n\sqrt r$, so \emph{every} critical point with $\theta>0$ breaks it once padded far enough: $S^4\times S^4$ already at $n=10$, $0.809040$ against $1.145644$. The bound factors as $\theta\le(n-1)|\lambda_{\min}|$, which holds at every critical point, composed with $|\lambda_{\min}|\le\sqrt{r/(N(N-r))}$, a claim of linear algebra with no curvature in it, and all the failure is in the second. The first is the $\lambda$-weighted bound of Lemma~41, whose equality class is already characterised: the slack vanishes exactly when $\Delta=0$ and $g$ is carried by the $\lambda_{\min}$ eigenspace. \emph{A successor states what its bound does when $W$ has a kernel, since $\dim E_+$ then measures the support rather than the space.}

\item \textbf{No per-row condition closes three or more distinct eigenvalues.} The complete per-row condition is a convex-hull membership (the row point $R_i=\big(\tfrac{-\nu_i}{n-2},\tfrac{\nu_i(\theta-\nu_i)}{n-2}\big)$ lying in the hull of the pair-points a multiplier uses), and at $n=12$ the spectrum with multiplicities $(30,20,16)$ satisfies every per-row condition while exceeding what the route needs. \emph{So the remaining content of the bracket-tensor route is cross-row.}

\item \textbf{The per-row multiplier calculus is a plane geometry, and it cannot close $d=3$.} A row multiplier is an affine $\ell(X,Y)$ against six pair-points ($P_{jj}=(\nu_j,\nu_j^2)$ on the parabola $Y=X^2$ and $P_{jk}=(\tfrac{\nu_j+\nu_k}2,\nu_j\nu_k)$ strictly below it), annihilating exactly the types on its line, with right-hand side $(n-2)\ell(R_i)$. Only a line through at least two of the six can forbid, and counting them settles the apparatus: of the $\binom62=15$ two-point lines the three tangents consume $9$, leaving \textbf{six per row and $27$ in all}. A tangent forbids a three-subset, lying in $7$ of the $\binom{10}4=210$ candidate supports, and a two-point line forbids one, so at most $9\times7+18=81$ die and \textbf{at least $129$ survive}. \emph{The argument is a count of the identities themselves and uses no enumeration of what they cut.}

\item \textbf{The half-flat single-four-set family carries exactly one critical point, so no exclusion theorem can come from it.} On that family $P$ is an $SO(3)$-invariant cubic on traceless symmetric $3\times3$ matrices, hence $\theta=9\det V_+$ for the self-dual block $V_+$. The gradient of $\det$ on the sphere forces two equal eigenvalues, so $\pm W_{\CP^2}$ is the only member, and the second-order condition can be evaluated nowhere else in it. \emph{A successor must clear a family with more than one critical point in it.}

\end{enumerate}

The isolation route's second link carries eight closures, and they are one closure with eight faces. \textbf{The pair clause is an equality wherever the orbit vector is nonzero} ($\widetilde N o=(\theta+M_{jk})o$ at every pure critical point), so no lossy estimate can certify it there, and orbit-free pairs are $45$--$78\%$ of all pairs rather than a thin exceptional set. \textbf{The inequality $\lambda_{\max}(\widetilde N)\le\theta+M_{jk}$ is false at padded pairs}, where the bound is exactly tight, and it was never what the clause asked: the clause is the separation $\theta+M_{jk}\notin\operatorname{spec}(\widetilde N|_{\mathbf1^\perp})$ off the orbit direction, which survives padding. \textbf{Conditions of the form $\langle(\widetilde N-T)y,x\rangle=0$ are inert} for every $y$, since $\ker\perp\operatorname{ran}$ makes them hold on the eigenspace by construction. \textbf{And every rearrangement of the pair-clause data is that eigen-equation restated}: the four-set discrepancy as orbit-vector differences, the extension of the orbit vector to all of $[n]$, the constant-residual reading, the approximate-eigenvector reading. Two checks catch such a rearrangement before a cycle is spent: whether the new form uses a quantity the old did not, and whether it beats the incumbent at one recorded point.

The four-set clause's own closures are three. \textbf{Concentration is refuted in both masses}: $m_4\ge\tfrac12$ does not force $\theta<\sqrt{3/2}$ ($S^5\times S^{43}$ is in band at $\theta=1.2636$ with $m_4=0.5012$), and neither does the row-mass reading, since four rows there carry $73\%$ of the total. \textbf{Both projection repairs fail in opposite ways}: onto $\mathcal D_n$ the result has support on all of $[n]$, so only $\Theta_n$ applies; onto the four-supported slice the extremal concentrator is orthogonal to a $2$-plane and the projection is identically zero. \textbf{And the complement coupling is determined by $\theta$ and four within-set row sums}, $X=\theta\sigma-\tfrac12\sum_{a\in S}\rho_a^2$, so reaching outside the four-set through it imports nothing.

Four posts bound what a successor may attempt on the operator itself. \textbf{The spectral gap has no $\Lambda^2$ analogue}: $\operatorname{tr}W=0$ and $\|W\|=1$ give $0<\lambda_{\max}<1<\theta$, so every unit Weyl operator in the band has an eigenvalue in $(0,\theta)$, criticality not needed. Lemma~129's gap is a statement about the $n\times n$ sectional matrix, and the two spectra are different objects, the $\Lambda^2$ one being the multiset of the sectional matrix's \emph{entries} on the diagonal locus. \textbf{And re-bounding $H=\mathrm{Ric}^*(\mathrm{Ric}(|W|))$ by trace data cannot close $J_5\ge0$}: any such bound is blind to $H$ scalar, where the cross term vanishes exactly, and carries a factor $\Theta(n^2)$. The recorded Cauchy--Schwarz is already the joint one, with $H$ inside its own coefficients, and the gap-commutator certificate, which vanishes on the commuting class, dominates it.

The remaining two close the $\#$-calculus rewrite of that same grant. \textbf{The three-term form $\langle X_{\mathrm{inv}},|W|\rangle+\tr(X_{\mathrm{inv}}\#|W|)-3\langle X_{\mathrm{inv}},\mathrm{Alt}(|W|)\rangle$ has zero $\Lambda^4$ content}: $\#I$ acts by $2$ on $\Lambda^4$, so the third term is exactly the counter-term and cancels it identically, leaving no defect information to extract beside $\tr D=6\|\mathrm{Alt}\|^2$; and the $\#$-trace collapses against $W$ or $I$ by $W\#I=-W$ and $I\#I=(n-2)I$, so $\tr(X_{\mathrm{inv}}\#W_-)$ is the one pairing carrying content. \textbf{The polarised scalar-curvature identity returns $J_5$ at every argument built from $W$'s spectral parts}, and the reason is a rank count: $\mathrm{Ric}(W)=0$ with $\mathrm{Ric}(W_\pm)=g$ makes $\mathrm{Ric}$ rank one on $\operatorname{span}\{W_+,W_-\}$, so evaluating the identity against $X_{\mathrm{inv}}$ at $B=\alpha W_++\beta W_-$ gives $\tfrac{\alpha+\beta}4J_5$ and nothing else. The $|W|$ and $W_-$ readings are therefore one instance, and $B=W$ gives zero, which is $W\#I=-W$ again. \emph{So the grant survives as a question about signed bracket weight ($\tr(X_{\mathrm{inv}}\#W_-)=\tfrac12\sum x_i\mu_j\|[u_i,v_j]\|^2$ with $x_i$ indefinite), and that is one of three unsigned terms rather than the whole of it.}

Four further posts close the stability-constant reading of Conjecture~L16 (\S\ref{sec:stability}), and three of them concern the competitor rather than the problem. \textbf{A partition argument cannot replace the two-regime comparison}: $P$ is additive only across orthogonal-support blocks, and a block sum of Weyl tensors is Weyl only on the granted diagonal locus, while the sphere and both champion orbits are compact so no scale parameter exists. \textbf{The $\ell_{a,b}$ transformation law is the identity on Weyl tensors}: $\mathrm{Ric}(W)=0$ gives $\ell_{a,b}(W)=W$ and $D_{a,b}(W)=0$ for every $(a,b)$, and on the spectral parts every term is a polynomial in $\mathrm{Ric}(W_\pm)$ alone, which the rank-one post already closes.

\textbf{At the competitor, no bound on $\delta^2$ alone reaches its second-order coefficient by any of three named routes.} A ceiling $\delta^2\le C$ with $C>\tfrac19$ leaves slack at the base point that absorbs any coefficient; global concavity of $\delta$ is impossible, $\delta$ being nonnegative, $1$-homogeneous and even, and the line-restricted version fails wherever it would say anything new, since a concave function cannot grow quadratically while $\delta^2(u+tH)/t^2\to\delta^2(H)$; and $u$ is not a local maximum of $\delta^2$ on the first-order equality cone, the zero-padded flat carrying $\delta^2=\tfrac3{25}>\tfrac19$ inside it. \emph{What defeats every attempt here is a feature of the competitor: it has a kernel of dimension $\binom n2-3$, so the zero-eigenvalue locus passes through it and the second derivative of $\delta^2$ carries a distributional atom that no pointwise identity can see.}

\emph{At the balanced orbit, by contrast, no such ceiling holds.} A tube of radius $\Theta(n^{-2})$ follows from the second-order margin and a Taylor remainder, but the exponent is not forced: the margin is attained on a pair block, where the defect's response vanishes identically, so the resolvent estimate that bounds the remainder is inert along the very direction that binds. What sets the exponent there is instead the crude bound $\sup\lvert P\rvert\le n-1$, and every critical point on record has $\sup\lvert P\rvert=O(1)$.

Fence post F34 has a positive counterpart worth stating, since it is a closed form the literature lacks: a candidate sharp bound on $\langle W^{\#},W\rangle$ exists and is exactly saturated by the balanced sphere product at every even $n$ tested.

\begin{Conj}{14}[Bound on $\langle W^{\#},W\rangle$]
For $W$ diagonal in coordinate $2$-planes, $\langle W^{\#},W\rangle=\tfrac12\tr(\Lambda^3)$ for the vertex matrix $\Lambda$, and Ricci-flatness forces $\Lambda\mathbf1=0$. The sharp bound on $\mathbf1^\perp$ then gives
\[
\langle W^{\#},W\rangle\;\le\;\sqrt2\cdot\frac{n-3}{\sqrt{(n-1)(n-2)}}\,\|W\|^3 ,
\]
saturated by the balanced sphere product at every even $n$ from $6$ to $20$, with the extremal spectral profile matching the theoretical one to six digits at $n=100$. \emph{What is unproven is the general case}: the identity is a fact about that frame rather than a tensorial one, and it does not survive rotation as stated.
\end{Conj}

\subsection{Three closures on the block induction}\label{sec:blockfences}

\paragraph{Freeing two outer blocks at fixed norms, in arbitrary diagonal fields, fails (post F68).} The fixed-norms form of the two-block step is false as stated over arbitrary fields $A$ and $C$: at $A=\diag(\sqrt2,-\sqrt2,0)$ and $C=\tfrac5{\sqrt2}\diag(0,1,-1)$ with $\lvert U\rvert=\lvert V\rvert=1$ an explicit witness beats it. \emph{The mechanism is majorisation}: for fixed $V$ the best $U$ makes a convex invariant, hence Schur-convex, function of $\tfrac94V^2+A$, largest when $V$'s squares are ordered like $A$'s diagonal. \textbf{What it does not touch is the block problem itself}, where the fields are built from the blocks rather than arbitrary, so a successor must use how they are built.

\paragraph{A condition on the Lagrange fields can be criticality restated (post F69).} A field defined by the criticality equation it came from yields conditions that are that equation again, and substituting the field's own form is the test that detects it. \textbf{Five instances are on record}, each reading as a second-order input and each vacuous on substitution. \emph{The target it appeared to supply, $[\hat U,C]=0$ at every global maximum, is untouched and still open.}

\subsection{Three closures around the four-set corner}\label{sec:fourfences}

\paragraph{No local-to-global theorem for forms on spheres applies to $P$ (post F73).} Padded $W_{\mathbb{CP}^2}$ is a spurious local maximum: its second variation is negative semi-definite on $W^{\perp}$ with excess $0$, which makes it a genuine local maximum, while $\theta=\sqrt{3/2}<\theta_0(n)$ for $n\ge12$ makes it spurious. \textbf{The higher-order conditions are vacuous there for the same reason}, excess $0$ making the third-order condition $T(e,e,e)=0$ and the fourth-order inequality carry no content. \emph{The exclusion is for arguments over the whole space}: one confined to a stratum the competitor is not in is untouched, which is what the four-set work below is.

\paragraph{The diagonal ceiling on $P(D)$ reaches $n\le10$ and misses $n=11$ (post F74).} Criticality forces $P(D)$ to an exact value on the four-set subalgebra, a function of two parameters on the equal-norm locus, while $D$ lies in $\mathcal D_n$ where the diagonal theorem caps $P$ at $\Theta_n\|D\|^3$, a cap unavailable for $W$ itself. \textbf{A configuration dies when the forced value exceeds the cap}, so what governs is the \emph{minimum} of the forced ratio over the admissible range, floored by the higher of two conditions rather than either alone. The shares excluded run $65.6, 56.5, 41.4, 30.8, 13.7$ per cent at $n=6,\dots,10$ and zero at $n\ge11$. \emph{Conjecture~A is proved at every $n\ge12$, so $n=11$ is the case the route must reach}, and the supremum of the forced ratio there is $1.21559$ against $\Theta_{11}=1.21716$, short by $0.13\%$.

\paragraph{Counting equations against unknowns cannot exclude the symmetric shape (post F75).} The scalar shadow of the diagonal equation is the sum of three of the six block equations rather than a seventh, given unit norm and the vanishing of $\langle D\circ D,N\rangle$; and $\theta-2L_{q_a}$ is invertible on the equal-projection locus, so each $\nu_a$ is determined by the other two and the system is square. \textbf{The component along $\nu_a$ is the pinning identity and already fixes that block's angle}, so counting the transverse components against the same unknowns double-spends them: the honest count is six equations against five unknowns. \emph{An exclusion at the symmetric shape must therefore come from the size of the solution the system has}, which is what closes it: the attainable $\tau$ there is capped at $1/\sqrt3$ of its bound against a floor of $0.5943$, so no admissible configuration exists. Two independent derivations, on identical briefs.

\paragraph{The $x\to\kappa_i$ link is real, sharp, and too small (post F70).} The identity $p+p'=\tfrac13\sum_{a<b\in Q_1}D_{ab}$ holds at every $D\in\mathcal D_n$, and in the span both $p$ and $p'$ are pinned, so the link exists and the sign of $x$ is not the obstacle. \textbf{It is a factor $1.65$ short}: at the surviving point it supplies $60.8\%$ of the $0.0148$ deficit, with every factor already at its own extreme, which over-estimates the joint maximum because the factors are coupled. Two independent audits, each rebuilding the ceiling from the definitions before folding it.

\subsection{Four closures at three eigenvalues}\label{sec:pairfences}

\paragraph{The class-level summed pair test cannot close three eigenvalues (post F76).} A relaxation witness at $n=12$, with multiplicities $(4,59,3)$, $e$-values $(-\tfrac17,0,\tfrac4{21})$ and $\theta=\tfrac{\sqrt{21}}2$, satisfies the whole three-value class system in exact rational arithmetic: the first-order rows, the pair caps, the Gram interval, and Lemma~235's two inequalities. \textbf{Modulo the rows, the summed test adds one moment}, and linear class-level data with the pair cap do not bound $\theta$ once a kernel is allowed, which the witness itself shows. The Ky Fan caps of Lemma~237 exclude the witness point but not its spectral ray, which still reaches $\tfrac{83\sqrt{21}}{294}\approx1.294$ exactly. \emph{A successor must carry an ingredient the ray violates}: a sharper cap for the three-plus-three partner profile (below $\approx7.724$ closes the ray at $\sqrt{3/2}$), a cap on a seven-frame's bracket total, or the unsummed per-pair inequalities, where Lemma~234's non-locality lives.

\paragraph{Each eigenspace tested on its own leaves a window of stars open from $n=19$ (post F77).} On the symmetric family $(-a,0,a)$ with $a=1/\sqrt{2p}$, let the negative eigenspace be a star $e\wedge U$ with $\dim U=p$. The per-vector system of Lemma~244 is feasible at $\theta=(p-3)/\sqrt{2p}$, which exceeds $\theta_0(n)$ once $p\ge9$, and Lemma~245 admits the star when $p\le\tfrac{n-1}2$. \textbf{So for $9\le p\le\tfrac{n-1}2$, non-empty from $n=19$ and growing linearly in $n$, a star passes both single-eigenspace tests}, and the window is empty at $12\le n\le18$. \emph{A successor must use the eigenspaces jointly}, which Lemma~246 does for stars; near-stars remain.

\paragraph{The traced rank certificate excludes no symmetric three-value critical point (post F78).} Lemma~247 fires wherever $\theta\sqrt{2p}>n-5$, and the traced rows with nonnegative traces force $\theta\sqrt{2p}\le n-5$ at every point they admit, so on the symmetric family the region where it fires holds no critical point. \emph{A certificate for the symmetric bulk must fire at $\theta\sqrt{2p}\le n-5$}, which needs an upper bound on $D=\tfrac12(t_{222}-t_{002}+p)$ that the traced rows do not supply. Off the symmetric family the certificate survives, and only where $\lvert\lambda_0\rvert>\lambda_2$.

\paragraph{The Bianchi identity, traced against the class four-forms, is rows already held (post F79).} On the symmetric family the first Bianchi identity reads $\Omega_0=\Omega_2$ in $\Lambda^4$, with $\Omega_c=\sum_{\psi\in E_c}\psi\wedge\psi$. Its three pairings $\langle\Omega_0-\Omega_2,\Omega_c\rangle=0$ are exactly three of the nine traced rows, and $\lVert\Omega_0-\Omega_2\rVert^2$ is their combination, so \textbf{traced Bianchi adds nothing to the rows}. \emph{A successor must use Bianchi pointwise}, as four-forms contracted against vectors or restricted to subspaces, or through class-level forms beyond the $\Omega_c$, each of which brings unknowns outside the rows.

\subsection{Two closures past three eigenvalues}\label{sec:latefences}

\paragraph{Uniqueness of positive-definite $\#$-idempotents is false (post F80).} At $SU(3)\subset G_2$, $C=\tfrac9{28}P_{\su(3)}+\tfrac3{28}P_{\pp}$ is a positive-definite idempotent on $\mathfrak g_2$ other than $I/4$ (Lemma~263), and Lemma~267 gives one on $\so(n)$ itself for every $n\ge7$. \textbf{So no argument deriving $C=I/c$ from $C\#C=C$ and $C\succ0$ alone can hold}, and neither can any pairing of the idempotent equation that uses nothing else. \emph{A successor must use more than positivity}: the second-order condition $C\#\cdot\le\tfrac12$ on $C^\perp$ that a maximiser's idempotent satisfies (Lemma~264), or the norm itself, and it must reach the idempotents lighter than $I/c$ (Lemma~270).

\paragraph{Small diagonal mass is separated from the maximiser only by the maximal frame (post F81).} Rotate the balanced product of $S^m\times S^m$, $n=2m$, by $\tfrac\pi4$ in the $m$ planes pairing its two factors. The result is critical with value $\theta_0(n)$, has no pair part, and has diagonal mass exactly $t^2=\tfrac1n$. \textbf{So at $n\ge12$ no strict bound at small diagonal mass follows from criticality and the sign-character sectors alone.} The witness frame is a critical point of the diagonal mass with negative semidefinite Hessian, and no single-plane rotation raises its diagonal mass, while rotating all matched planes together does, at fourth order. \emph{A successor must use the maximal frame at fourth order across several planes at once.}

% ==============================================================================
